I need to analyze this sparse signal detection problem, which relates to the Donoho-Jin detection boundary theory.

Given parameters:
- beta = 0.85 (sparsity)
- r = 0.45 (signal strength)
- mu_n = sqrt(2r log n)
- epsilon_n = n^(-beta)

Let me search for the specific detection boundaries for these tests.Based on my research, I can now analyze this sparse signal detection problem using the Donoho-Jin detection boundary theory.

**Given parameters:**
- beta = 0.85 (sparsity parameter)
- r = 0.45 (signal strength parameter)  
- mu_n = sqrt(2r log n)
- epsilon_n = n^(-beta)

**The detection boundary:**

For sparse normal means testing with beta > 1/2, the detection boundary separates the beta-r plane into the detectable region and the undetectable region. The optimal detection boundary is rho*(beta) = beta - 1/2 for beta in (1/2, 1).

With beta = 0.85: rho*(0.85) = 0.85 - 0.5 = 0.35

Since r = 0.45 > 0.35 = rho*(beta), we are in the **detectable region**.

**Analysis of each test:**

**A) One-sided z-test based on Xbar:**
The z-test based on the sample mean works when signals are dense (beta <= 1/2). If the sparsity level is such that alpha <= 1/2, then ANOVA is asymptotically optimal. Since beta = 0.85 > 1/2, the signal is too sparse for the z-test to have any power. The expected signal in the mean is n^(1/2-beta) x sqrt(2r log n) -> 0, so the z-test has power -> alpha, **NOT power -> 1**.

**B) Max test:** 
Check if $\sqrt{r} + \sqrt{1-\beta} > 1$: $\sqrt{0.45} + \sqrt{0.15} = 0.671 + 0.387 = 1.058 > 1$. Thus, following the problem's guidance, the max test doesn't achieve full power here.

**C) Higher Criticism test:**
Donoho and Jin (2004), following work of Ingster (1999), studied the problem of testing for any signal in a sparse normal means model and showed that there is a "detection boundary" above which the signal can be detected. They showed that Tukey's "higher criticism" statistic achieves the detection boundary.

The higher criticism is successful throughout the same region of amplitude sparsity where the likelihood ratio test would succeed. Since it does not require a specification of the alternative, this shows that higher criticism is in a sense optimally adaptive to unknown sparsity and size of the nonnull effects.

Since r = 0.45 > rho*(0.85) = 0.35, the HC test **has power -> 1**.

**Answer: [C]**
